\begin{table*}[t]
\centering
\small
\caption{\textbf{Examples of application (A) levels.}
One question from CLIcK per level, with our English translation and the reason for its label.
The gold answer is in \textbf{bold}.}
\label{tab:a_examples}
\begin{tabular}{@{}p{0.09\textwidth}p{0.29\textwidth}p{0.31\textwidth}p{0.23\textwidth}@{}}
\toprule
 & \textbf{Question (Korean)} & \textbf{English translation} & \textbf{Why this level} \\
\midrule
A0 \newline Knowledge retrieval
& 대한민국 헌법재판소에는 몇 명의 헌법재판관이 있는가?
\newline A) 7명 \quad B) 8명 \quad \textbf{C) 9명} \quad \mbox{D) 10명}
& How many justices sit on the Constitutional Court of Korea?
\newline A) 7 \quad B) 8 \quad \textbf{C) 9} \quad D) 10
& The recalled fact is itself the answer; there is no case to map it onto. \\
\midrule
A1 \newline Explicit mapping
& 이혼 당사자 간의 협의가 잘 이루어지지 않을 때, 자녀의 양육권은 어떻게 결정되는가?
\newline \textbf{A) 가정법원에서 결정한다}
\newline B) 자녀의 의사에 따라 결정한다
\newline C) 변호사가 결정한다
\newline D) 이혼한 당사자 간의 협의에 의해 결정한다
& When the divorcing parties cannot reach an agreement, how is custody of the children decided?
\newline \textbf{A) The family court decides}
\newline B) It is decided by the wishes of the children
\newline C) A lawyer decides
\newline D) It is decided by agreement between the divorced parties
& The question states the condition of the rule (the parties fail to agree), so the rule maps onto the case directly. \\
\midrule
A2 \newline Inferred mapping
& 예금액이 많을 경우 어떻게 해야 모두 보호받을 수 있는가?
\newline A) 한 은행에 모두 예금한다
\newline \textbf{B) 여러 금융 기관에 나누어 예금한다}
\newline C) 금액을 줄인다
\newline D) 보호받을 수 없다
& If one has a large amount of deposits, what should one do so that all of it is protected?
\newline A) Deposit all of it in one bank
\newline \textbf{B) Divide it among several financial institutions}
\newline C) Reduce the amount
\newline D) It cannot be protected
& One rule applies (deposits are protected up to a limit per institution), but it is not stated; one must infer that a large deposit exceeds the limit and what would satisfy the rule. \\
\midrule
A3 \newline Mapping selection
& 다음 정치관과 관련이 깊은 정책으로 옳은 것은?
\newline 임금의 직책은 한 사람의 재상을 논정하는 데 있다 하였으니, 바로 총재를 두고 한 말이다. 총재는 위로는 임금을 받들고 밑으로는 백관을 통솔하여 만민을 다스리는 것이니 직책이 매우 크다. (후략)
\newline A) 육조 직계제의 시행
\newline B) 사간원의 독립
\newline \textbf{C) 의정부 서사제의 시행}
\newline D) 집현전의 설치
& Which policy is closely related to the following view of government?
\newline ``It is said that the duty of the king lies in choosing one chancellor. The chancellor serves the king above and leads all officials below in governing the people, so his office is of great weight.'' (abridged)
\newline A) Direct reporting of the Six Ministries to the king
\newline B) Independence of the Office of the Censor-General
\newline \textbf{C) Deliberation by the State Council}
\newline D) Establishment of the Hall of Worthies
& Two institutions both govern how the king and his ministers share power, with opposite effects: direct reporting strengthens the king, and deliberation by the State Council strengthens the chancellors. One must select which of the two the stated view corresponds to. \\
\bottomrule
\end{tabular}
\end{table*}